% Generated draft: requires mathematical and source review.
\documentclass[11pt]{article}
\usepackage[margin=25mm]{geometry}
\usepackage{fontspec}
\setmainfont{Latin Modern Roman}
\usepackage{amsmath,amssymb,mathtools,unicode-math}
\setmathfont{Latin Modern Math}
\usepackage{enumitem,needspace,xurl,hyperref}
\setlength{\parindent}{0pt}
\setlength{\parskip}{5pt}
\setlength{\emergencystretch}{3em}
\newcommand{\N}{\mathbb{N}}
\newcommand{\Z}{\mathbb{Z}}
\newcommand{\Q}{\mathbb{Q}}
\newcommand{\R}{\mathbb{R}}
\newcommand{\C}{\mathbb{C}}
\newcommand{\F}{\mathbb{F}}
\newcommand{\A}{\mathbb{A}}
\newcommand{\PP}{\mathbb{P}}
\newcommand{\E}{\mathbb{E}}
\newcommand{\eps}{\varepsilon}
\newcommand{\dd}{\,\mathrm{d}}
\begin{document}
\section*{UnsolvedMath ID: 8700031}
\textbf{Source problem code:} catalogue entry
\subsection{Definitions and mathematical statement}
Consider a family of polynomial inputs $P_1,\ldots,P_m$ in a common, variable number of variables $x_1,\ldots,x_n$, represented by a total of $N$ coefficients. The promise is that the polynomials have no common zero in $\mathbb C^n$. The task is to output coefficients of polynomials $A_1,\ldots,A_m$ satisfying the polynomial identity
\[
\sum_{i=1}^m A_iP_i=1.
\]
The coefficient field, input monomial representation, and output representation are not specified. Interpreting the polynomials in $\mathbb C[x_1,\ldots,x_n]$ gives one natural algebraic setting, but the exact computational field must be fixed with the model.

The conjecture excludes a polynomial-step algorithm: for every fixed absolute exponent $c>0$, no single program solves this search problem on every promised input in at most $N^c$ steps. The exponent is independent of the number of variables, number of polynomials, degrees, and input coefficients. The source's attribution places the question in the Blum--Cucker--Shub--Smale complexity context, but the allowed arithmetic/test operations and constants are not supplied.

An exact formulation must distinguish arithmetic-operation complexity from bit complexity, dense from sparse coefficient encodings, and explicit coefficient output from straight-line/circuit output. The conjectured absence of every polynomial-step bound is not by itself a specified exponential lower bound, despite the research summary's wording. It is a search-complexity question, not an automatically established NP-hardness assertion.
\subsection{Short English statement}
Given polynomials with no common complex zero, prove that no polynomial-step program always outputs explicit B\'ezout identity coefficients. Fix the coefficient encoding, output representation, and computational model before stating the lower bound.
\subsection{Sources}
\begin{itemize}
\item Michel Waldschmidt, \emph{Open Diophantine Problems} (2004), Conjecture 3.13, \url{https://arxiv.org/abs/math/0312440}; accessed 2026-07-29, source-TeX extraction. The item is attributed to Blum, Cucker, Shub, and Smale.
\item Local registry \texttt{data/unsolvedmath/problems.json}, permanent UnsolvedMath ID $8700031$, source code catalogue entry, \emph{Conjecture 3.13 --- Blum, Cucker, Shub and Smale}. Dataset attribution: Ulam AI and the UnsolvedMath Contributors.
\item The research summary connects the question to BCSS complexity and the effective Nullstellensatz but supplies no separate primary publication or model specification.
\end{itemize}
\end{document}
